% ============================================================
% ［架空の記入例・提出不可］結果，解釈および今後の課題はすべて架空である．
% 通常のchapters/06-conclusion.texには自動で読み込まれない．
% ============================================================
\chapter{\UsamiConclusionTitle}

\begin{center}
  \UsamiFictionalExampleNotice
\end{center}

\section{本研究のまとめ}

本研究では，低照度の単眼RGB動画から得られた2次元人物骨格を対象として，誤推定骨格を後段処理の前に除外する方法を検討した．既存の姿勢推定器は再学習せず，人物ごとの平均関節信頼度と，胴体長で正規化した骨格辺長の基準からの逸脱を統合して骨格品質スコアを定義した．さらに，動画系列単位で学習・検証・評価を分け，検証系列で決定した閾値を固定して評価系列へ適用した．

架空の低照度条件では，提案スコアにより採用骨格のPCK@0.2が除外なしより8.9ポイント，平均関節信頼度だけによる除外より3.9ポイント高くなった．一方，採用率は除外なしより18.3ポイント低下した．この結果は，構造的に不自然な骨格を優先的に除外できた可能性を示すが，姿勢推定器自体の精度向上を意味しない．また，精度の向上と引き換えに骨格の欠損が増えるため，後段処理が許容できる採用率を考慮して閾値を選ぶ必要がある．

以上より，本例の範囲では，関節信頼度と骨格構造を併用することが出力骨格の品質管理に有用である可能性を確認した．ただし，使用した動画，数値および結論は組版と論述のための架空例であり，実在する研究成果ではない．

\section{今後の課題}

第1に，時間情報を利用していないため，単一フレームでは低品質でも前後フレームから補間可能な骨格を除外する場合がある．骨格軌跡の速度と加速度を用いた判定を追加し，誤除外と見逃しの変化を確認する必要がある．第2に，基準辺長を単一の撮影環境から求めており，体格，撮影距離および視点の違いに対する頑健性は未評価である．人物ごとの尺度推定と複数環境での外部評価が必要となる．第3に，採用率だけでなく，動作分類や人数計測等の後段タスクの性能を直接測り，用途ごとに閾値を決定する必要がある．

実研究へ展開する場合は，撮影対象者の同意とデータ管理手順を明示し，照度，遮蔽，動きおよび人物属性ごとの失敗例を報告する．これらを確認した上で，単純な後処理として導入できる範囲と，再学習が必要な範囲を切り分けることが課題である．
